%! Functions and Function Notation — Unit 1, Lesson 1.1 — Student Packet Cover Sheet
% The ONE place \namedateperiod belongs (Namestrip). 1.1 now OPENS the unit, so nothing here
% refers back to a previous lesson.
\documentclass[10pt]{article}
\usepackage{precalculus-article}
\usepackage{precalculus-boxes}
\usepackage{ltablex}
\keepXColumns

\begin{document}

% Full-bleed plum banner
\begin{tikzpicture}[remember picture, overlay]
  \fill[plum] (current page.north west)
    rectangle ([yshift=-0.9in]current page.north east);
\end{tikzpicture}
\vspace{-0.8in}

\begin{center}
  {\color{white}\textbf{\LARGE Precalculus}}\\[4pt]
  {\color{lilacmid}\large\textbf{Unit 1: Functions and Change}}\\[6pt]
  {\color{white}\textbf{\large Lesson 1.1 \quad Functions and Function Notation}}
\end{center}

\vspace{0.22in}
\namedateperiod
\vspace{0.18in}

\begin{learningtargetbox}
\small
\begin{enumerate}[leftmargin=1.5em, itemsep=5pt]
  \item I can read and write function notation, and name the three facts a
        statement like $f(3) = 7$ carries: the input, the output, and a point on
        the graph.
  \item I can evaluate a function given by a formula, a table, or a graph ---
        including at negative inputs.
  \item I can state the domain and range of a function given as a list of
        points or as a graph.
  \item I can evaluate a \textit{piecewise} function by deciding which rule owns
        the input --- including when the input sits exactly on a boundary.
\end{enumerate}
\end{learningtargetbox}

\vspace{0.14in}

\begin{tocbox}
\small
\renewcommand{\arraystretch}{1.5}
\begin{tabularx}{\linewidth}{@{} c l X r @{}}
\rowcolor{lilac}
\textbf{\#} & \textbf{Component} & \textbf{Description} & \textbf{Score} \\
\midrule
1 & Warm-Up      & Spiral review of prerequisite skills            & \blank{1.2cm} \\
2 & Guided Notes & Notation, evaluating, domain and range, piecewise functions & \blank{1.2cm} \\
\rowcolor{goldbg}
3 & AP Practice  & \textbf{Extra credit} --- AP-style multiple choice and free response & $+$\,\blank{1.0cm} \\
4 & Homework     & Graded practice in new contexts \quad Due: \blank{2.2cm} & \blank{1.2cm} \\
\midrule
  & \multicolumn{2}{r}{\textbf{Total} \quad {\footnotesize (1, 2, and 4, plus any extra credit)}} & \blank{1.2cm} \\
\end{tabularx}
\end{tocbox}

\vspace{0.14in}

\begin{remindbox}
\small
A function is a machine you can trust: put in an input, get back exactly one
output. Today that machine gets a \textbf{name} --- $f$ --- so one short symbol
can carry the rule, the input, and the output all at once. Two words describe
everything it can take in and put out: \textbf{domain} and \textbf{range}. And
one function may be built from several rules at once --- a \textbf{piecewise}
function --- where each rule owns its own stretch of inputs.
\end{remindbox}

\end{document}
